% ---------------------------------------------------------------------------
% Chapter 6: The Family Photograph: Generator-Level Jets and Truth Definitions
% Generated from the outline; edit freely, this file is now the source.
% ---------------------------------------------------------------------------
\chapter{The Family Photograph: Generator-Level Jets and Truth Definitions}
\chaptermark{The Family Photograph}
\label{ch:family_photograph}

\begin{journeybox}
Before the family leaves we take a photograph. Everyone who later claims to be a
member will be compared against it. But even a photograph involves choices: do
we include the neutrinos, and who decides whether this is a $b$ family at all?
\end{journeybox}

\physicsline{GenJets, jet flavour definitions (ghost association, hadron vs parton flavour, flavour-$k_T$, IFN), truth for JEC}

\section{What we mean by ``the true jet'': choices and their consequences}
\label{sec:family_photograph:what_we_mean_by_the_true_jet_choices_and}
\todo{Write this section.}

\section{Ghost-associated hadron flavour vs parton flavour}
\label{sec:family_photograph:ghost_associated_hadron_flavour_vs_parto}
\todo{Write this section.}

\section{Infrared-safe flavour: flavour-$k_T$ and interleaved flavour neutralisation}
\label{sec:family_photograph:infrared_safe_flavour_flavour_k_t_and_in}
\todo{Write this section.}

\section{Our jet at this stage: the reference GenJet we will compare to forever}
\label{sec:family_photograph:our_jet_at_this_stage_the_reference_genj}
\todo{Numbers, plots and event display of the $b$-jet at this stage.}

\begin{ledger}{The Family Photograph}
  \ledgerrow{before this stage}{}{}{}{--}
  \ledgerrow{after this stage}{}{}{}{}
\end{ledger}

\section{Further reading}
\label{sec:family_photograph:further_reading}
Official and theory references for this stage: \cite{Banfi:2006hf}, \cite{Caletti:2022hnc}, \cite{Caletti:2022glq}, \cite{Cacciari:2008gp}.

\begin{pitfalls}
\todo{Pitfalls and FAQs for newcomers at this stage.}
\end{pitfalls}

\begin{exercises}
\item \todo{Exercise 1.}
\item \todo{Exercise 2.}
\end{exercises}
